\section{The antidiagonal local classes}\label{sec:local}

Set $L_+=L$, $L_-=K_C\otimes L^{-1}$ and
$Y=\Tot(L_+\oplus L_-)$. The torus acts trivially on $C$ and by
$(v_+,v_-)\mapsto(tv_+,t^{-1}v_-)$ on the fibers. Denote by $Z_m$
the space of compact effective one-cycles in $Y$ of degree $m$ over
$C$. Every irreducible component maps finitely onto $C$. Indeed, an
affine plane fiber contains no complete curve, and the projection
of each compact component is proper. The space $Z_0$ is the point
representing the empty cycle. Since $K_Y$ is trivial, all these
cycles have Chern number zero and belong to the semi-Fano locus.

\subsection{Symmetric functions on the fibers}
For $p\in C$ the fiber intersection of $D\in Z_m$ is a multiset of
$m$ points in $Y_p$. Project it to each of $L_{+,p}$ and $L_{-,p}$.
For $1\leq j\leq m$ the $j$th elementary symmetric function of a
projection gives an element
\[
 b_{\pm,j}(D)(p)\in L_{\pm,p}^{\otimes j}.
\]
These evaluations are continuous functions of $D$ by the
fiber-intersection map of \cite[(4.15)]{Pardon}.

\begin{lemma}\label{lem:support}
For $m>0$, choose $1+\max\{0,m(g-1)\}$ distinct points of $C$.
The common zero set of the functions $b_{\pm,j}$ evaluated at these
points consists of $m[C]$. This cycle is also the unique fixed point
of the circle on $Z_m$.
\end{lemma}
\begin{proof}
The coefficients $b_{\pm,j}(D)$ are holomorphic sections on $C$.
To see this across a branch point, normalize each irreducible
component. Its finite map to the smooth base curve is flat. In a
local trivialization of the line, the fiber coordinate acts by
multiplication on the finite locally free direct image of the
structure sheaf. Its characteristic polynomial records the projected
branches with multiplicities. Take the product of these polynomials
over the components, raising each to its cycle multiplicity. The
$j$th coefficient, up to the sign $(-1)^j$, transforms as a
section of $L_\pm^{\otimes j}$.

Both line bundles have degree $g-1$. If all the chosen evaluations
vanish, each coefficient vanishes identically by the degree bound.
The two projected polynomials are thus $T^m$ on every fiber. All
points of the multiset have both coordinates zero, which forces
$D=m[C]$. Conversely $m[C]$ is a common zero. For a circle-fixed
cycle the circle preserves each finite fiber multiset. Connectedness
forces it to fix each point individually. The origin is the only
fixed point of the fiber action.
\end{proof}

For later use, $Z_m$ is finite dimensional. Embed $Y$ in the bundle
of lines in $\cO_C\oplus L_+\oplus L_-$. On the open chart $Y$
the relative $\cO(1)$ is trivial. Tensoring it with a sufficiently
ample line bundle pulled back from $C$ gives a projective
polarization whose degree on $D$ is $m$ times its degree on the
zero section. Thus $Z_m$ is the open subset of a fixed-degree
projective Chow space parametrizing cycles disjoint from the divisor
at infinity. This description gives the compact-cycle topology on
that open set and the bounded dimension needed for equivariant
localization.

\subsection{The finite Borel construction}
Write $B_N=\PP^N$, $h=c_1(\cO_{B_N}(1))$, and
\[
 Z_{m,N}=\big((\CC^{N+1}\setminus\{0\})\times Z_m\big)/\CC^\times
        \longrightarrow B_N
\]
for the diagonal quotient. It is the relative degree-$m$ cycle
space of the corresponding Borel family of local threefolds.
Trivializing $L_\pm$ at each chosen point makes the tests of
Lemma~\ref{lem:support} into functions of weights $+j$ and $-j$.
A function $f$ of weight $r$ gives a section of the pullback of
$\cO_{B_N}(r)$. Our sign convention follows from the quotient:
the class of $(v,s)$ of weight one defines the functional on
$\CC v$ taking $v$ to $s$.

For any tuple $F=(f_1,\ldots,f_n)$ of equivariant functions of
nonzero weights $r_i$ with compact common zero set, put
\begin{equation}\label{eq:cut}
 u_N(F)=\frac{1}{\prod_i r_i}\,
    F^*\operatorname{Th}\!\left(\bigoplus_{i=1}^{n}\cO_{B_N}(r_i)\right)
        \quad\text{in } H_c^{2n}(Z_{m,N};\QQ).
\end{equation}
Here the Thom class is taken relative to the complement of the
zero section. Its pullback has support in the common zero set over
$B_N$, which is compact. In particular the definition allows
noncompact individual zero sets and needs no transversality. On
forgetting supports the normalized Euler class of a single section
is $c_1(\cO(r))/r=h$ for either sign of $r$.

\begin{lemma}\label{lem:stabilization}
The universal image of $u_N(F)$ depends only on $n-N$ for the given
local model and degree.
\end{lemma}
\begin{proof}
For the inclusion $i:B_N\hookrightarrow B_{N+1}$ of a hyperplane
use the same symbol for the inclusion of Borel families.
The class $u_{N+1}(F)$ restricts to $u_N(F)$, so the hyperplane
Gysin map gives
\[
 i_!u_N(F)=h\,u_{N+1}(F)=u_{N+1}(F,f)
\]
for every additional function $f$ of nonzero weight. The second
equality uses the compact support already supplied by $F$, which
allows the last relative Euler class to be replaced by its ordinary
class. The first equality is the projection formula for the complex
normal line pulled back from the base. The product description
locally over the base suffices even for a singular cycle space.

These Gysin maps induce the stabilization of relative classes in
\cite[(4.5)--(4.7)]{Pardon}. Given two tuples of lengths $n,n'$ at
stages $N,N'$ with $n-N=n'-N'$, append the second tuple to the
first and the first to the second. Both products now live at
stage $N+n'=N'+n$, and the even Thom factors commute. Hence the
images in the universal group coincide.
\end{proof}

Take $n=N$ in \eqref{eq:cut} and let $a_m(C,L)$ denote the image
in Pardon's algebra. Such a tuple exists by Lemma~\ref{lem:support}.
The oriented base has real dimension $2N$, so a compact class of
degree $2n$ has relative degree $2n-2N$ \cite[Lemma~3.5]{Pardon}.
The class $a_m(C,L)$ has virtual dimension and Chern number zero
and therefore lies in $H$. In particular its expression in the
scalar generators is a polynomial. The counit evaluates on the
empty cycle, hence vanishes on $a_m(C,L)$ for $m>0$. We set
$a_0(C,L)=1$, as given by the empty tuple at $N=0$.
